% Audited research proof; stronger bounds appear in the following refinement.
% Include after random-host-weighted-matching.tex and random-hosts-phase5.tex.
\section{Keeping wheels with partially useful length sets}
\label{sec:rhost-one-fifty-one}

A wheel remains useful when some of its cycle lengths exceed the target interval.
Retain its useful lengths and charge its full rank.
This avoids the main loss in the preceding construction.

\begin{lemma}[Partial wheel length sets]\label{lem:rhost-partial-wheels}
Let $h\to\infty$ through multiples of nine.
There is a list of traceable wheel subdivisions with marked cycle-length sets $A$ such that:
\begin{enumerate}
\item each $A$ contains its rim length in $[3,h/3+1]$ and at least four non-rim lengths in $[h+1,2h]$;
\item the marked sets are pairwise disjoint;
\item $\sum_A(|A|-4)\ge(35/48-o(1))h$.
\end{enumerate}
Every wheel satisfies the same core thresholds and edge bounds as Lemma~\ref{lem:rhost-wheel-profit}.
\end{lemma}

\begin{proof}
Put $s=\lceil\sqrt h\rceil$ and use the previous spoke distribution.
Choose $R$ uniformly from $\{6s,\ldots,\lfloor h/12\rfloor\}$.
Choose $L$ uniformly from the integer interval with center $\lfloor(R-2)/2\rfloor$ and radius $\lfloor s/2\rfloor$.
Put $M=R-2-L$ and use the rim $(1,L,1,M)$.
Its marked rim length is $R$.
Retain only those non-rim cycle lengths that lie in $[h+1,2h]$.
Let $k$ be their number, counted initially by cycle labels.
The prospective value of this outcome is $k-3$.

Each spoke-pair deficit is uniform on $\{0,\ldots,h+U\}$.
For a rim arc of length $a$, its cycle exceeds $2h$ with probability at most $(2s+a)/h$.
The two arcs for each of the six pairs sum to $R$.
All non-rim cycles have length at least $h+1$.
Consequently,
\[
 \mathbb E(12-k)\le\frac{24s+6\mathbb E R}{h}
 =\frac14+o(1),
\]
and
\[
 \mathbb E(k-3)\ge\frac{35}{4}-o(1).
\]
Discard outcomes with a nonpositive spoke or repeated non-rim cycle lengths.
Their total probability is $o(1)$, by the calculations in the preceding section.
The value is bounded, so this changes its expectation by $o(1)$.
Also discard outcomes with $k\le3$.
Their values are nonpositive, so this cannot lower the retained expected value.
Each retained marked set has size $k+1\in\{5,\ldots,13\}$ and value $|A|-4\in\{1,\ldots,9\}$.

Every upper-interval vertex has weighted degree at most $12/h$.
Every rim vertex has weighted degree at most
\[
 \frac1{\lfloor h/12\rfloor-6s+1}=\frac{12+o(1)}h.
\]
The weighted codegree bound $O(h^{-3/2})$ is unchanged.
Restricting the marked set to its useful cycle lengths only decreases these degrees and codegrees.
For a rim and a non-rim length, fixing $R$ has probability $O(1/h)$ and fixing the pair deficit costs at most $1/h$.
Apply the variable-size version of Lemma~\ref{lem:rhost-weighted-matching}, with maximum value nine.
It supplies disjoint marked sets of total value at least
\[
 \frac{h}{12+o(1)}\left(\frac{35}{4}-o(1)\right)
 =\left(\frac{35}{48}-o(1)\right)h.
\]
Choose one metric witnessing each selected set.
The core bounds, total edge count, and spanning path follow exactly as in Lemma~\ref{lem:rhost-wheel-profit}.
\end{proof}

\begin{theorem}[An upper coefficient of $151/144$]\label{thm:rhost-one-fifty-one}
For the same $p_*$ as in Theorem~\ref{thm:rhost-scale}, and every $p\ge p_*$,
with high probability
\[
 \operatorname{Pex}(G(n,p))
 \le\log_2n+\left(\frac{151}{144}+o(1)\right)\frac{\ln n}{\ln\ln n}.
\]
\end{theorem}

\begin{proof}
Use the short-block baseline of Theorem~\ref{thm:rhost-seventy-nine}, of rank $16h/9+O(1)$.
For every marked set in Lemma~\ref{lem:rhost-partial-wheels}, replace its $|A|$ single-cycle blocks by the corresponding wheel.
This saves $|A|-4$ in cycle rank.
The disjointness of the marked sets makes these replacements compatible.
The resulting rank is at most
\[
 \frac{16h}{9}-\left(\frac{35}{48}-o(1)\right)h
 =\left(\frac{151}{144}+o(1)\right)h.
\]
The remaining $o(d)$ single-cycle blocks, simultaneous random embeddings, connector chain, and later exposures are unchanged.
All discarded cycle lengths belong to an embedded wheel; they need not belong to its marked set.
Every required length still has either its retained baseline block or a marked replacement.
The previous excess calculation proves the assertion.
\end{proof}
